\chapter{User guide}
\label{chap:user-guide}

This chapter explains how to use FlowSync, role by role, as a short practical manual. It follows
the pages of the frontend (Chapter~\ref{chap:functionalities}); the interface is in Italian, so
the names of buttons and pages are given in Italian, with a translation where useful.

\status{All the steps below can be followed today in the frontend running against its mock
backend (Section~\ref{sec:dev-run}): the demonstration data are those of a small workshop with
nine resources and four orders. Plans, indicators, alerts and audit records are produced by the
mock; with the real backend, the same steps will work as its modules are completed. Registration,
login, resources, catalogues and orders already exist in the real backend.}

\section{Getting started}
\label{sec:ug-start}

Open the address of the application (in development \texttt{http://localhost:5173}) and log in
on the page \emph{Accedi}. Three test accounts exist, one per role, both in the mock and in the
development database; their password is \texttt{demo}, and with the mock the fields can also be
left empty to enter as a planner.

\begin{table}[htpb]
  \centering
  \small
  \begin{tabular}{@{}lll@{}}
    \toprule
    \textbf{Username} & \textbf{Role} & \textbf{First page} \\
    \midrule
    \texttt{pianificatore} & planner & dashboard \\
    \texttt{operatore} & operator & activities of the day \\
    \texttt{amministratore} & administrator & dashboard (with Users and Audit in the menu) \\
    \bottomrule
  \end{tabular}
  \caption{Test accounts (development only)}
  \label{tab:test-accounts}
\end{table}

A new person asks for an account from the link \emph{Registrati} on the login page, giving name,
username, e-mail and a password of at least eight characters. The account is created
\emph{pending}: it can be used only after an administrator has approved it and chosen its role
(Section~\ref{sec:ug-admin}). The name, the role and the button \emph{Esci} (log out) are at the
bottom of the side menu, or at the top of the page for operators.

\section{Planner}
\label{sec:ug-planner}

\paragraph{1. Prepare the orders.} In \emph{Commesse} (orders) the list on the left shows the open
orders; \emph{Nuova commessa} creates one, \emph{Modifica dati} changes customer, dates, priority
and weight of the delay. The due date may be left empty: the system will ask for it before
planning. Below, the graph shows the activities of the selected order:
\begin{itemize}
  \item \emph{Aggiungi attività} adds an activity: name, required skills and one or more
        alternatives, each with one resource and its duration in hours; tick
        \emph{Durata non nota} when the duration is unknown, and the AI will estimate it;
  \item drag from the dot on the right of an activity to the dot on the left of another to say
        that the second starts after the first; click an arrow to remove it. The column
        \emph{Dopo} (after) of the table does the same from the keyboard;
  \item nothing is stored until \emph{Salva modifiche}; \emph{Annulla modifiche} restores the
        saved version.
\end{itemize}

\paragraph{2. Ask for the plans.} In \emph{Nuova pianificazione} write the request in words,
for example ``plan C-101 to C-104 for next week, avoiding overtime'' (\emph{Usa un esempio}
fills in a sample text), and press \emph{Interpreta}. Check what the system understood: orders, period, hard constraints and preferences,
which become the three weights (punctuality, overtime, setups); move the sliders if needed and
answer the questions about missing data. Then press \emph{Genera i 3 piani}: the page shows the
progress and, after a few seconds, the result.

\paragraph{3. Choose a plan.} The comparison page shows three plans: \emph{Economico}
(slower, cheaper), \emph{Bilanciato} (balanced, the recommended one) and \emph{Veloce} (faster,
more expensive), each with total duration, tardiness, overtime, setups, utilisation and cost,
and a table with the same indicators for the ATC and EDD rules. \emph{Apri nel Gantt} opens a
plan.

\paragraph{4. Review and decide.} On the Gantt chart each row is a resource and each bar an
activity. Drag a bar to move it in time or to another row; if the move is not allowed, a message
explains why and the bar goes back. Selecting a bar opens a panel with the details and with
buttons to move it by 15 minutes or an hour, or to another resource. Then:
\begin{itemize}
  \item \emph{Approva} makes the plan the running one and discards the other two alternatives;
  \item \emph{Rifiuta} rejects it, with a reason;
  \item after some moves, \emph{Salva come nuova versione} stores the modified plan, with a
        reason, as a new plan to approve.
\end{itemize}
All the requests and plans are listed in \emph{Piani}, with search and filters.

\paragraph{5. When something goes wrong.} When an operator reports a problem, an alert appears
on the dashboard. Its link \emph{ripianifica} opens the rescheduling page, which already knows
the problem: check the activities affected, optionally set until when the machine will be down,
the weights and how much the plan may change, and press \emph{Ripianifica con l'AI}. The result
is a new request with three plans, to be compared and approved as in steps 3 and 4; the running
plan is replaced only when a new one is approved. The button \emph{Ripianifica} on the Gantt
chart of an approved plan does the same, asking first what happened if there is no report.

\section{Operator}
\label{sec:ug-operator}

The page \emph{Le tue attività di oggi} lists the activities of the day. Press \emph{Inizia}
(start) when you begin an activity, \emph{Pausa} and \emph{Riprendi} to interrupt and resume it,
\emph{Termina} when it is finished. Only one activity can be in progress, and an activity cannot
start before the previous ones of its order are finished: the page says what you are waiting for.
To report a problem, choose its kind in \emph{Segnala un problema} (machine breakdown, missing
material, delay, other), the machine involved and, if useful, a note, and send it: the planner
receives it at once.

\section{Administrator}
\label{sec:ug-admin}

\paragraph{Users.} The page \emph{Utenti} shows first the pending registrations (their number is
also next to the menu entry): choose the role and press \emph{Approva}, or \emph{Rifiuta} to
delete the request. Below, the list of the accounts allows changing a role or disabling and
re-enabling an account. Your own account cannot be disabled or lose the administrator role.

\paragraph{Configuration.} In \emph{Configurazione} four tabs hold the master data: resources
(type, skills, calendar, hourly cost, planned unavailability such as a maintenance), calendars
(working days, shifts, closures), skills and domain profiles (default weights and active
constraints). These data are used by every planning request: for instance, a maintenance window
entered here is respected by every new plan.

\paragraph{Audit and model.} \emph{Audit} lists who did what and when, with filters and the
export \emph{Esporta CSV}; each record shows the model version and the hashes that make the plan
reproducible. \emph{Modello AI} shows which model versions are active and how often the fallback
rules were used.

\section{Troubleshooting}
\label{sec:ug-troubleshooting}

\begin{itemize}
  \item \emph{``Il tuo account aspetta ancora l'approvazione''}: the account is pending; an
        administrator must approve it.
  \item \emph{``Sessione scaduta''} or the login page appearing by itself: the session has
        expired; log in again.
  \item A section shows \emph{``Il servizio non è disponibile''} with a retry button: the backend
        did not answer. With the mock, this message can be produced on purpose (see
        Section~\ref{sec:dev-frontend}).
  \item A move on the Gantt chart is refused: the message names the constraint, typically an
        activity that must finish first, a resource already busy, or a resource not allowed for
        that activity.
\end{itemize}
